\documentclass[11pt]{article}
\usepackage[margin=1in]{geometry}
\usepackage[T1]{fontenc}
\usepackage{lmodern,amsmath,amssymb,amsthm,mathtools}
\usepackage[colorlinks=true,allcolors=blue]{hyperref}
\newtheorem{proposition}{Proposition}
\newtheorem{corollary}[proposition]{Corollary}
\newcommand{\R}{\mathbb R}
\newcommand{\dd}{\,\mathrm d}
\newcommand{\norm}[1]{\left\lVert#1\right\rVert}
\newcommand{\curl}{\operatorname{curl}}
\title{Separating forcing and low-frequency feeding\protect\\in a signed Navier--Stokes tail budget}
\author{}
\date{September 27, 2026}
\begin{document}
\maketitle
\begin{abstract}
For a smooth finite-energy forced Navier--Stokes evolution, we bound the
absolute time integrals of direct force work and low-frequency feeding
in the high-frequency enstrophy balance. The bounds use kinetic energy,
total energy dissipation, a fixed cutoff, and derivatives of the force;
they spend none of the high-frequency viscous dissipation. Consequently,
for the forced singularities constructed in~\cite{OAI}, the accumulated
tail self-stretching minus even the full tail dissipation is unbounded
at every fixed cutoff. The full-field obstruction is checked in Lean.
The Fourier-tail extension below is an analytic proof, not a formalized
projector theorem. We also bound the forcing derivative of a quartic
caloric stretching functional using only energy-level velocity controls
and smooth force norms. Its signed nonlinear renewal must therefore fail
a uniform absorption budget on the same forced construction. A Fourier
multiplier estimate separates every low-frequency parent and every low
renewal output. Low stretching outputs have a separate energy-paid
history bound, leaving a target with every parent and both internal
frequencies high. That target is unproved; no unforced regularity
estimate is obtained.
\end{abstract}

\section{An exact split and finite allowances}

Let $u,p,f$ solve
\[
 \partial_tu+(u\cdot\nabla)u=-\nabla p+\nu\Delta u+f,
 \qquad \nabla\cdot u=0,\qquad \nu>0,
\]
on $\R^3\times[a,T)$, where $0<L=T-a<\infty$. Assume smoothness and, on
each compact slab $[a,b]\subset[a,T)$, $u\in C^1([a,b];H^m(\R^3))$
and $f,\nabla p\in C([a,b];H^m(\R^3))$ for every integer $m\geq0$.
These hypotheses justify the following strong energy identities and
integrations by parts. Write
\[
 E(t)=\norm{u(t)}_2^2,\quad Y(t)=\norm{\nabla u(t)}_2^2
       =\norm{\omega(t)}_2^2,\quad \omega=\curl u.
\]
Suppose
\[
 E(t)\leq M,\quad \norm{f(t)}_2\leq F_0,
 \quad\norm{\Delta f(t)}_2\leq F_\Delta\qquad(a\leq t<T).
\]
All constants in this note are independent of the shorter endpoint $t$.
The kinetic-energy identity gives
\begin{equation}\label{eq:B}
 \int_a^tY(s)\dd s\leq B,
 \qquad B=\frac{E(a)}{2\nu}+\frac{L\sqrt M F_0}{\nu}.
\end{equation}
Indeed, $E'/2=-\nu Y+\langle u,f\rangle$ and $E(t)\geq0$.

Fix an angular Fourier cutoff $K\geq0$, and set
\[
 u_L=\mathbf1_{\{|\xi|<K\}}(D)u,\quad u_H=u-u_L,
 \quad\omega_H=\curl u_H,\quad S_H=\tfrac12(\nabla u_H+\nabla u_H^\mathsf T).
\]
These are orthogonal $L^2$ projections commuting with derivatives.
Define the physical tail quantities
\[
 Z_H=\tfrac12\norm{\omega_H}_2^2,\qquad
 D_H=\norm{\nabla\omega_H}_2^2=\norm{\Delta u_H}_2^2,
 \qquad \Pi_H=\int_{\R^3}\omega_H\cdot S_H\omega_H\dd x.
\]
In particular, $\Pi_H$ uses the high field in all three factors. With
$\Pi=\int\omega\cdot S\omega$, introduce
\[
 P_L=\langle-\Delta u_L,u\times\omega\rangle,
 \quad P_H=\langle-\Delta u_H,u\times\omega\rangle,
 \quad X_H=P_H-\Pi_H,\quad F_H=\langle-\Delta u_H,f\rangle.
\]
Then $P_L+P_H=\Pi$, and the exact tail balance is
\begin{equation}\label{eq:balance}
 Z_H'=\Pi_H+X_H-\nu D_H+F_H.
\end{equation}
The filtered velocity is not assumed to solve an autonomous
Navier--Stokes equation: $X_H$ retains every lower-frequency input.

\begin{proposition}\label{prop:allowances}
Let $c_K=(K^3/(6\pi^2))^{1/2}$. Under the preceding assumptions,
\begin{align}
 \int_a^t|F_H(s)|\dd s&\leq L\sqrt M F_\Delta,\label{eq:force}\\
 \int_a^t|X_H(s)|\dd s&\leq
  3Kc_K\sqrt M B+K^2c_KM\sqrt{LB}=:C_L.\label{eq:low}
\end{align}
Neither bound uses a gradient supremum or tail dissipation of the velocity.
\end{proposition}
\begin{proof}
Integration by parts gives $F_H=-\langle u_H,\Delta f\rangle$, so
$|F_H|\leq\sqrt M F_\Delta$. This proves~\eqref{eq:force}.

For the unitary Fourier convention, Cauchy--Schwarz on the ball gives
\[
 \norm{\nabla u_L}_\infty\leq Kc_K\sqrt{E},\qquad
 \norm{\Delta u_L}_\infty\leq K^2c_K\sqrt{E}.
\]
Here the gradient uses the Frobenius norm, and
$c_K^2=(2\pi)^{-3}|B_K|$. The vector and matrix bounds follow by
Cauchy--Schwarz before taking a component sum.

For matrices define the trilinear form
\[
 b(A,B,C)=\sum_{i,j,k} A_{ji}B_{kj}C_{ki},
 \qquad |b(A,B,C)|\leq |A|_F|B|_F|C|_F.
\]
The integrated stretching identity is
$\Pi=-\int b(\nabla u,\nabla u,\nabla u)$, and the same identity holds
for the smooth solenoidal field $u_H$. Put $A=\nabla u$, $H=\nabla u_H$
and $G=\nabla u_L=A-H$. The exact telescoping identity
\[
 b(A,A,A)-b(H,H,H)=b(G,A,A)+b(H,G,A)+b(H,H,G)
\]
and $\norm{H}_2\leq\norm{A}_2$ imply
\begin{equation}\label{eq:cubic}
 |\Pi-\Pi_H|\leq 3\norm{\nabla u_L}_\infty Y
       \leq3Kc_K\sqrt M\,Y.
\end{equation}
Also
\[
 |P_L|\leq\norm{\Delta u_L}_\infty\norm{u}_2\norm{\omega}_2
       \leq K^2c_K M\sqrt Y.
\]
Now $X_H=\Pi-\Pi_H-P_L$. Integrate these two bounds and use
\eqref{eq:B} and $\int_a^t\sqrt Y\leq\sqrt{LB}$.
\end{proof}

\begin{corollary}\label{cor:gradient}
If a finite constant $A$ satisfies
\begin{equation}\label{eq:target}
 \int_a^t(\Pi_H-\nu D_H)\dd s\leq A\qquad(a\leq t<T),
\end{equation}
then the complete gradient energy obeys
\[
 Y(t)\leq 2Z_H(a)+2A+2C_L+2L\sqrt M F_\Delta+K^2M.
\]
In particular, a half-dissipation budget implies this conclusion, but
the strictly smaller fraction is unnecessary for this fixed-cutoff
gradient-energy implication.
\end{corollary}
\begin{proof}
Integrate~\eqref{eq:balance}, apply Proposition~\ref{prop:allowances},
and use $Y=2Z_H+\norm{\nabla u_L}_2^2\leq2Z_H+K^2M$.
\end{proof}

For an unforced solution, one can take $M=E(a)$ and $B=E(a)/(2\nu)$.
Thus the low-frequency allowance is controlled by the datum, viscosity,
cutoff and horizon. This does not prove~\eqref{eq:target}: its uniform
signed high-frequency history remains the open estimate. Allowing the
full dissipation simply exposes its relation to the gradient-energy
continuation criterion; it is not a new coercive estimate.

\section{Application to the actual forced construction}

The whole-space construction in~\cite{OAI} has smooth compactly supported
forcing, bounded kinetic energy, zero initial velocity, and unbounded
speed as $t\uparrow1$. On slabs exhausting $[1/2,1)$, all hypotheses
above hold with fixed $M,F_0,F_\Delta$.

The module \texttt{MettapediaNS.\allowbreak ForcedSignedBudget} in the Lean
library applies directly to the construction's \texttt{CandidateProperties}.
It proves
that its gradient energy cannot have a uniform bound: such a bound
would control the order-two energy via the forced viscous estimate,
and the $H^2$ embedding would bound speed. It also proves the full-field
signed-production obstruction directly, with force work bounded by
$6\sqrt{E_0(u)}\sqrt{E_2(f)}$ in the doubled energy convention, where
$E_j(v)=\norm{\nabla^jv}_2^2$ sums all ordered derivative components.
The constructor theorem is imported at the pinned revision
\texttt{f9e8bc5b38b6e212696e8a30e3e91517af887bbd}.

\begin{corollary}\label{cor:forced}
For every such candidate, every fixed finite $K\geq0$, and every
$\beta\leq1$,
\[
 \sup_{1/2\leq t<1}\int_{1/2}^t
       (\Pi_H(s)-\beta\nu D_H(s))\dd s=+\infty.
\]
Meanwhile $\int_{1/2}^1|X_H|$ and $\int_{1/2}^1|F_H|$ are finite.
\end{corollary}
\begin{proof}
For $\beta=1$, a finite upper bound would give a uniform gradient-energy
bound by Corollary~\ref{cor:gradient}, contradicting the preceding
forced-candidate theorem. For $\beta\leq1$, the accumulated quantity
only increases because $D_H\geq0$. The other assertions follow from
Proposition~\ref{prop:allowances}.
\end{proof}

Thus the signed tail self-interaction is the term that defeats the
budget; an infinite direct forcing-work allowance is not the cause.
This conclusion applies to the complete constructed solution, without
replacing its oscillatory corrections by a finite numerical model.

Forcing is still essential to that constructed history. With its zero
datum and zero force, the kinetic-energy identity forces the smooth
finite-energy solution to remain zero. In the construction the force
is the smoothly extended momentum residual of the designed field
(Section~10 of~\cite{OAI}). Removing it while keeping that field would
change the equation. Restarting an unforced evolution from a later
nonzero slice is a different problem, whose subsequent history is not
settled by this argument. A proposed unforced budget must therefore
use more than bounded forcing norms or bounded direct force work.

\section{Forcing in a caloric stretching functional}

The preceding obstruction can also test a functional that records
stretching beyond the instantaneous enstrophy. Let
\[
 A=-\Delta,\qquad H_s=e^{-\nu sA},\qquad
 R=\curl A^{-1}P,\qquad
 \mathcal M(v)=R\bigl((\curl v\cdot\nabla)v\bigr),
\]
where $P$ is the whole-space Leray projection. The multiplier $R$ is
used on a divergence: $\mathcal M(v)=R\operatorname{div}(v\otimes\curl v)$.
Thus its low-frequency singularity causes no $L^2$ difficulty. Parseval's
identity gives $\norm{R\operatorname{div}Q}_2\leq\norm Q_2$, where the
tensor norm is Frobenius. For solenoidal $v$,
\begin{equation}\label{eq:Mbound}
 \norm{\mathcal M(v)}_2\leq\norm v_\infty\norm{\nabla v}_2,
 \qquad \langle Av,\mathcal M(v)\rangle=\Pi(v).
\end{equation}
Define the nonnegative functional
\begin{equation}\label{eq:Kdef}
 \mathcal K(u)=\frac1\nu\int_0^\infty
                  \norm{\mathcal M(H_su)}_2^2\dd s.
\end{equation}
It is finite for the smooth $L^2$ fields considered here: the integrand
is bounded near zero and is $O(s^{-5/2})$ at infinity by heat smoothing.

The effective force in the projected velocity equation is $g=Pf$, not
necessarily $f$ itself. Suppose in addition that
\[
 \norm{g(t)}_\infty\leq Q_0,\qquad
 \norm{\curl g(t)}_\infty\leq Q_1\qquad(a\leq t<T).
\]
These constants are finite for the smooth compactly supported force in
the construction: $P$ is bounded on every $H^m$, Sobolev embedding controls
$Pf$, and $\curl Pf=\curl f$. No $L^\infty$ boundedness of $P$ is assumed.

\begin{proposition}\label{prop:caloricforce}
Put $c_h=(8\pi)^{-3/4}$. For a fixed smooth solenoidal slice $u$ and
direction $g$ as above, with $E=\norm u_2^2$ and $Y=\norm{\nabla u}_2^2$,
\begin{equation}\label{eq:caloricforce}
 |D\mathcal K(u)[g]|\leq\frac{c_h}{\nu^2}
 \left(\frac{32}{3}Q_0 E^{3/4}Y^{3/4}
                 +16Q_1 E^{5/4}Y^{1/4}\right).
\end{equation}
Consequently the standing evolution satisfies, uniformly for $t<T$,
\begin{equation}\label{eq:caloricforceintegral}
 \int_a^t|D\mathcal K(u)[g]|\dd s\leq C_{\mathcal K}:=
 \frac{c_h}{\nu^2}\left(
 \frac{32}{3}Q_0 M^{3/4}L^{1/4}B^{3/4}
       +16Q_1 M^{5/4}L^{3/4}B^{1/4}\right).
\end{equation}
This payment spends no order-two velocity dissipation.
\end{proposition}
\begin{proof}
Use the auxiliary heat parameter $r=\nu s$, and write
$w_r=e^{-rA}u$, $h_r=e^{-rA}g$. Differentiating both parent slots gives
\[
 D\mathcal M(w)[h]
 =R\operatorname{div}(h\otimes\curl w+w\otimes\curl h).
\]
The multiplier bound, heat contraction and~\eqref{eq:Mbound} imply
\begin{align*}
 |D\mathcal K(u)[g]|\leq\frac2{\nu^2}
 \int_0^\infty \norm{w_r}_\infty\norm{\nabla w_r}_2
 \bigl(Q_0\norm{\nabla w_r}_2+Q_1\sqrt E\bigr)\dd r.
\end{align*}
The Gaussian kernel has $L^2$ norm $c_hr^{-3/4}$, and its Fourier
multiplier gives
\[
 \norm{w_r}_\infty\leq c_hr^{-3/4}\sqrt E,\qquad
 \norm{\nabla w_r}_2\leq\min\{\sqrt Y,r^{-1/2}\sqrt E\}.
\]
When $E,Y>0$, split each integral at $r_*=E/Y$. Direct integration yields
\begin{align*}
 \int_0^\infty\norm{w_r}_\infty\norm{\nabla w_r}_2^2\dd r
 &\leq\frac{16}{3}c_h E^{3/4}Y^{3/4},\\
 \sqrt E\int_0^\infty\norm{w_r}_\infty\norm{\nabla w_r}_2\dd r
 &\leq 8c_h E^{5/4}Y^{1/4}.
\end{align*}
These bounds also justify differentiation under the heat integral for
the smooth slices. If $E$ or $Y$ vanishes, $u=0$ as an $L^2$ whole-space
field and the result is immediate. This proves~\eqref{eq:caloricforce}.
Finally apply H\"older in physical time, using $E\leq M$,
$\int Y\leq B$, and the interval length at most $L$.
\end{proof}

Let $N(u)=-P((u\cdot\nabla)u)$ be the actual unforced velocity source, and set
\[
 \mathcal G=\frac{Y+\mathcal K}{2},\qquad
 \mathcal T=D\mathcal K(u)[N(u)],\qquad
 \mathcal W=\nu\norm{Au}_2^2+
                   \frac1\nu\norm{\mathcal M(u)-\nu Au}_2^2.
\]
In particular $\mathcal G\geq Y/2$ and $\mathcal W\geq0$.
Heat translation in~\eqref{eq:Kdef} gives
$D\mathcal K(u)[-\nu Au]=-\norm{\mathcal M(u)}_2^2/\nu$.
Combine this with the actual forced enstrophy identity and expand the
square in $\mathcal W$ to obtain the exact balance
\begin{equation}\label{eq:caloricbalance}
 2\mathcal G'+\mathcal W
 =\mathcal T+D\mathcal K(u)[g]+2\langle Au,f\rangle.
\end{equation}
Both extra force terms have absolutely bounded physical-time integrals:
their sum is bounded by $C_{\mathcal K}+2L\sqrt M F_\Delta$.

\begin{corollary}\label{cor:caloricforced}
On every forced candidate of the preceding section, for every $\beta\leq1$,
\[
 \sup_{1/2\leq t<1}\int_{1/2}^t
          (\mathcal T-\beta\mathcal W)\dd s=+\infty.
\]
\end{corollary}
\begin{proof}
Integrating~\eqref{eq:caloricbalance} expresses
$\int(\mathcal T-\mathcal W)$ as
$2\mathcal G(t)-2\mathcal G(1/2)$ minus a uniformly bounded force integral.
The candidate's gradient energy has no uniform upper bound, and
$\mathcal G\geq Y/2$. This proves the case $\beta=1$; decreasing $\beta$
only increases the integral since $\mathcal W\geq0$.
\end{proof}

Thus the smooth force does not hide an infinite payment in the derivative
of this quartic functional. Its signed nonlinear renewal is itself
unbounded relative to its positive cost on the forced history. A proof
of the corresponding unforced budget must exploit the zero-force
evolution; estimating these explicit force terms is insufficient.
The estimate does not bound $\mathcal T$: substituting $N(u)$ for $g$
would require velocity-dependent supremum norms that are not controlled
by the energy inequality.

\section{Low parents and high-frequency caloric renewal}

The caloric return admits a further frequency separation. Fix $K\geq0$,
write $v=P_{\geq K}u$, $\ell=u-v$, and retain $E=\norm u_2^2$,
$Y=\norm{\nabla u}_2^2$, $D=\norm{\Delta u}_2^2$.
Choose a universal homogeneous Sobolev constant $S$ such that
$\norm w_6\leq S\norm{\nabla w}_2$ for real scalar, vector and matrix
fields on $\R^3$. The constant $c_K$ is as in Proposition~\ref{prop:allowances}.

\begin{proposition}\label{prop:lowparents}
If $E\leq M$, then for every $\varepsilon>0$,
\begin{equation}\label{eq:lowparents}
 |\mathcal T(u)-\mathcal T(v)|
 \leq\varepsilon\nu D+C_2Y^2+C_{3/2}Y^{3/2},
\end{equation}
where all constants depend only on $K,\nu,M,\varepsilon$ and $S$:
\[
 a_K=\frac{2c_KS^3}{\nu^2},\qquad
 b_K=\frac{6c_KS^3}{\nu^2},\qquad
 C_2=a_K\sqrt M+\frac{3b_K^{4/3}M}{4(4\varepsilon\nu)^{1/3}},
 \qquad C_{3/2}=a_KKM.
\]
In particular, the two last terms in~\eqref{eq:lowparents} are at most
$2(C_2Y+C_{3/2}\sqrt Y)\mathcal G$. Their coefficient is integrable
in physical time under the energy bounds.
\end{proposition}
\begin{proof}
We first bound the complete quartic Fourier kernel of $\mathcal K$.
For ordered inputs $a,b$ let
$\mathcal M_b(a,b)=R((\curl a\cdot\nabla)b)$.
At frequencies $p,q$ with output $k=p+q$, its coefficient satisfies
\begin{equation}\label{eq:orderedkernel}
 |\widehat{\mathcal M_b}(p,q)|
 \leq |p|\,|\widehat a(p)|\,|\widehat b(q)|.
\end{equation}
Indeed, $p\times\widehat a(p)$ is perpendicular to $p$, so its dot
product with $q$ equals its dot product with $k$. This cancels the
$|k|^{-1}$ multiplier of $R$. The zero output is zero as well. Thus no
lower bound on the output frequency is needed.

In an ordered quartic pairing of two such sources, heat integration gives
the denominator $\nu^2(|p|^2+|q|^2+|r|^2+|s|^2)$. Its kernel is bounded by
\begin{equation}\label{eq:quartickernel}
 \frac{|p|\,|r|}{\nu^2(|p|^2+|q|^2+|r|^2+|s|^2)}
 \leq\frac1{2\nu^2}.
\end{equation}
All parent pairs and all interference between their heat rates are included.

For a real field $a$, let $a^\#$ be the real scalar field with Fourier
transform $|\widehat a|$. It has the same $L^2$ norm and homogeneous
Sobolev norms as $a$. Its physical values need not be nonnegative.
The absolute ordered quartic pairing is bounded by $1/(2\nu^2)$ times
the nonnegative Fourier convolution corresponding to
$\int a^\#b^\#c^\#d^\#$. Applying H\"older and Sobolev after this
Fourier majorization gives
\begin{equation}\label{eq:derivativeL2}
 |D\mathcal K(u)[h]|\leq\frac{2}{\nu^2}
                    \norm{u^\#}_6^3\norm h_2
 \leq\frac{2S^3}{\nu^2}Y^{3/2}\norm h_2.
\end{equation}
There are four differentiated slots. Subtracting the derivatives at $u$
and $v$ telescopes each of the three other slots, giving twelve terms.
Since $v$ is a Fourier restriction of $u$, coefficient majorization yields
\begin{align}
 |(D\mathcal K(u)-D\mathcal K(v))[h]|
 &\leq\frac6{\nu^2}\norm{\ell^\#}_\infty
                         \norm{u^\#}_4^2\norm h_2\notag\\
 &\leq\frac{6c_KS^{3/2}}{\nu^2}
                      E^{3/4}Y^{3/4}\norm h_2.\label{eq:derivativeDifference}
\end{align}
The coefficient comparison takes place inside the Fourier convolution;
it does not assert a pointwise inequality between $|v^\#|$ and $|u^\#|$.
Here we used $\norm{\ell^\#}_\infty\leq c_K\sqrt E$ and interpolation
$\norm{u^\#}_4^2\leq S^{3/2}E^{1/4}Y^{3/4}$.

Write $\mathcal B(a,b)=-P((a\cdot\nabla)b)$. Then
\[
 N(u)-N(v)=\mathcal B(\ell,u)+\mathcal B(v,\ell),\qquad
 \norm{N(u)-N(v)}_2\leq c_K\sqrt{EY}+Kc_KE.
\]
Also Sobolev and interpolation give
\[
 \norm{N(v)}_2\leq\norm v_6\norm{\nabla v}_3
       \leq S^{3/2}Y^{3/4}D^{1/4}.
\]
Only the $L^2$ contraction and derivative commutation of the sharp
projector are used; no general $L^p$ bound for the ball multiplier is assumed.
Split the return exactly as
\[
 \mathcal T(u)-\mathcal T(v)
 =D\mathcal K(u)[N(u)-N(v)]
       +(D\mathcal K(u)-D\mathcal K(v))[N(v)].
\]
Equations~\eqref{eq:derivativeL2}--\eqref{eq:derivativeDifference} bound
these terms by
\[
 a_K(E^{1/2}Y^2+KEY^{3/2})
                  +b_KE^{3/4}Y^{3/2}D^{1/4}.
\]
Young's inequality with exponents $4$ and $4/3$ bounds the last term by
$\varepsilon\nu D+3b_K^{4/3}EY^2/(4(4\varepsilon\nu)^{1/3})$.
Use $E\leq M$ and $Y\leq2\mathcal G$ to finish.
\end{proof}

Low renewal outputs are paid separately, even when both source parents
are arbitrarily high. Set
\[
 g_K=P_{<K}N(v),\qquad
 \mathcal T_K(u)=D\mathcal K(v)[P_{\geq K}N(v)].
\]
The full source $N(v)$ is formed before splitting its output.
Cauchy--Schwarz in the unitary Fourier convolution gives
$|\widehat{N(v)}(\xi)|\leq(2\pi)^{-3/2}|\xi|E$.
Integrating this over the ball, with an extra $|\xi|$ for curl, proves
\begin{equation}\label{eq:lowoutputnorms}
 \norm{g_K}_\infty\leq\frac{K^4E}{8\pi^2},\qquad
 \norm{\curl g_K}_\infty\leq\frac{K^5E}{10\pi^2}.
\end{equation}
The derivative is the output frequency because the nonlinearity is in
divergence form; this estimate assumes no bound on parent gradients.
Proposition~\ref{prop:caloricforce}, applied at the slice $v$, therefore gives
\begin{equation}\label{eq:lowoutputpayment}
 \int_a^t|D\mathcal K(v)[g_K]|\dd s\leq C_{\rm out},
\end{equation}
where $C_{\rm out}$ is the right side of~\eqref{eq:caloricforceintegral}
with $Q_0=K^4M/(8\pi^2)$ and $Q_1=K^5M/(10\pi^2)$.
Although $v$ is not an autonomous solution, its energy and integrated
gradient energy are bounded by those of $u$, which is all that this
application uses.

\begin{corollary}\label{cor:highrenewal}
Suppose $A_0\geq0$ and $\theta<1$ satisfy, on every shorter prefix,
\begin{equation}\label{eq:highrenewaltarget}
 \int_a^t\mathcal T_K(u)\dd s
            \leq A_0+\theta\int_a^t\mathcal W(u)\dd s.
\end{equation}
Choose $0<\varepsilon<1-\theta$, and let
$C_f=C_{\mathcal K}+2L\sqrt M F_\Delta$ for the forced equation, or $C_f=0$
when $f=0$. Then
\[
 \mathcal G(t)\leq
 \left(\mathcal G(a)+\frac{A_0+C_{\rm out}+C_f}{2}\right)
           \exp\!\left(C_2B+C_{3/2}\sqrt{LB}\right).
\]
\end{corollary}
\begin{proof}
The exact split is
$\mathcal T(u)=\mathcal T_K(u)+\mathcal T(u)-\mathcal T(v)
                       +D\mathcal K(v)[g_K]$.
Insert this in~\eqref{eq:caloricbalance}, integrate, and use
Propositions~\ref{prop:caloricforce} and~\ref{prop:lowparents}.
Since $\nu D\leq\mathcal W$, the remaining cost has coefficient
$1-\theta-\varepsilon>0$ and can be dropped. Gr\"onwall applies with
coefficient $C_2Y+C_{3/2}\sqrt Y$, whose integral is at most
$C_2B+C_{3/2}\sqrt{LB}$.
\end{proof}

For an unforced regularity argument, $K$ and $A_0$ in
\eqref{eq:highrenewaltarget} must be chosen from the full initial datum,
viscosity and a prescribed finite horizon, uniformly over every shorter
lifespan and prefix. That high-frequency signed-history estimate is not
proved here. The result pays low parents and low outputs; it does not
pay repeated high-frequency regeneration.

The same estimates also sharpen the forced test: on every actual forced
candidate, for each fixed $K$ and every $\theta<1$,
\[
 \sup_{1/2\leq t<1}\int_{1/2}^t
       (\mathcal T_K(u)-\theta\mathcal W(u))\dd s=+\infty.
\]
A finite upper bound would give a uniform gradient-energy bound by
Corollary~\ref{cor:highrenewal}, contradicting the constructed candidate.
Unlike Corollary~\ref{cor:caloricforced}, this argument requires
$\theta<1$ because the low-parent estimate spends an arbitrarily small,
positive part of the dissipation.

\section{Low stretching outputs along the actual history}

The output of $\mathcal M$ is distinct from the output of the nonlinear
direction $N$. It too can be separated without restricting its parents.
Fix $J\geq0$, put $Q_-=P_{<J}$, $Q_+=P_{\geq J}$, and define
\[
 \mathcal K_\pm(u)=\frac1\nu\int_0^\infty
          \norm{Q_\pm\mathcal M(H_su)}_2^2\dd s,\qquad
 \mathcal Z_\pm(u)=\frac1\nu\norm{Q_\pm\mathcal M(u)}_2^2,\qquad
 \mathcal T_\pm(u)=D\mathcal K_\pm(u)[N(u)].
\]
Orthogonality gives $\mathcal K=\mathcal K_-+\mathcal K_+$ and
$\mathcal T=\mathcal T_-+\mathcal T_+$. All source outputs are formed
before the projection is applied.

\begin{proposition}\label{prop:lowstretching}
With $c_J^2=J^3/(6\pi^2)$,
\begin{equation}\label{eq:lowstretchingbounds}
 0\leq\mathcal K_-(u)\leq\frac{c_J^2E^2}{4\nu^2},\qquad
 0\leq\mathcal Z_-(u)\leq\frac{c_J^2EY}{\nu}.
\end{equation}
Along the actual forced evolution,
\begin{equation}\label{eq:lowstretchinghistory}
 \frac{\dd}{\dd t}\mathcal K_-(u)
       =-\mathcal Z_-(u)+\mathcal T_-(u)+D\mathcal K_-(u)[g].
\end{equation}
Consequently, uniformly for all shorter prefixes,
\[
 \left|\int_a^t\mathcal T_-(u)\dd s\right|
 \leq\frac{c_J^2M^2}{4\nu^2}+\frac{c_J^2MB}{\nu}+C_{\mathcal K}.
\]
The last term is zero for an unforced evolution. This bounds the signed
primitive, not the integral of $|\mathcal T_-|$.
\end{proposition}
\begin{proof}
The ordered bound~\eqref{eq:orderedkernel} and Cauchy--Schwarz in the
full convolution give, at every nonzero output $k$,
\[
 |\widehat{\mathcal M(u)}(k)|
 \leq(2\pi)^{-3/2}\int |p|\,|\widehat u(p)|\,
                  |\widehat u(k-p)|\dd p
 \leq(2\pi)^{-3/2}\sqrt{EY}.
\]
Integrating its square over $|k|<J$ proves the $\mathcal Z_-$ bound.
No upper bound on the parent frequencies was used. If $E_s=\norm{H_su}_2^2$
and $Y_s=\norm{\nabla H_su}_2^2$, then $E_s'=-2\nu Y_s$, so
\[
 \mathcal K_-(u)\leq\frac{c_J^2}{\nu}\int_0^\infty E_sY_s\dd s
                  \leq\frac{c_J^2E^2}{4\nu^2}.
\]
Heat translation gives $D\mathcal K_-(u)[-\nu Au]=-\mathcal Z_-(u)$.
The actual equation therefore proves~\eqref{eq:lowstretchinghistory}.
Proposition~\ref{prop:caloricforce} applies to $\mathcal K_-$ with
unchanged constants: inserting an orthogonal output projection in both
factors of its derivative can only improve the two $L^2$ norm bounds.
Integrate the identity, use
$|\mathcal K_-(u(t))-\mathcal K_-(u(a))|\leq c_J^2M^2/(4\nu^2)$,
and use $\int\mathcal Z_-\leq c_J^2MB/\nu$.
\end{proof}

For the unforced equation there is a sharper payment, independent even
of the horizon. The actual kinetic-energy equality is $E'=-2\nu Y$,
so~\eqref{eq:lowstretchingbounds} gives
\begin{equation}\label{eq:lowstretchingreserve}
 \mathcal K_-(u(t))+\int_a^t\mathcal Z_-(u)\dd s
 \leq \frac{c_J^2 E(a)^2}{4\nu^2}.
\end{equation}
Indeed, the bound for the integral is
$c_J^2(E(a)^2-E(t)^2)/(4\nu^2)$, which complements the bound for
the remaining capacity. Therefore every unforced prefix satisfies
\[
 -\mathcal K_-(u(a))\leq\int_a^t\mathcal T_-(u)\dd s
 \leq\frac{c_J^2 E(a)^2}{4\nu^2}-\mathcal K_-(u(a)).
\]
This is a bound on repeated low-output renewal from the initial kinetic
energy; its constant grows with the fixed output cutoff.

There is a more direct way to use this payment in the regularity target.
Keep the original, full positive cost $\mathcal W$, and set
$\mathcal G_+=(Y+\mathcal K_+)/2\geq Y/2$. Subtracting
\eqref{eq:lowstretchinghistory} from~\eqref{eq:caloricbalance} gives
the exact identity
\begin{equation}\label{eq:highstretchingbalance}
 2\mathcal G_+'+\mathcal W
 =\mathcal T_++\mathcal Z_-+D\mathcal K_+(u)[g]+2\langle Au,f\rangle.
\end{equation}
Thus only the spent low-output cost, bounded by
$C_{\rm src}=c_J^2MB/\nu$, is needed on the right. The endpoint
$\mathcal K_-$ allowance is unnecessary in this identity. In the unforced
case one may take $C_{\rm src}=c_J^2E(a)^2/(4\nu^2)$ by
\eqref{eq:lowstretchingreserve}.

With the input cutoff $K$ from the preceding section, define
\[
 v=P_{\geq K}u,\qquad
 \mathcal T_{K,J}(u)=D\mathcal K_+(v)[P_{\geq K}N(v)].
\]
The estimates of Proposition~\ref{prop:lowparents} hold with
$\mathcal K_+$ in place of $\mathcal K$: its ordered quartic kernel
is restricted at the common output, and has the same bound.
The force estimate and the low nonlinear-output payment hold as well.
Consequently, if fixed $A_0\geq0$ and $\theta<1$ obey
\begin{equation}\label{eq:fullyhightarget}
 \int_a^t\mathcal T_{K,J}(u)\dd s
          \leq A_0+\theta\int_a^t\mathcal W(u)\dd s
\end{equation}
on every shorter prefix, the same Gr\"onwall argument gives
\[
 \mathcal G_+(t)\leq
 \left(\mathcal G_+(a)+\frac{A_0+C_{\rm src}+C_{\rm out}+C_f}{2}\right)
       \exp\!\left(C_2B+C_{3/2}\sqrt{LB}\right).
\]
Here $0<\varepsilon<1-\theta$ defines $C_2,C_{3/2}$ as before. The
exact decomposition used in~\eqref{eq:highstretchingbalance} is
\begin{align*}
 \mathcal T_+(u)={}&\mathcal T_{K,J}(u)
       +D\mathcal K_+(u)[N(u)-N(v)]\\
 &+(D\mathcal K_+(u)-D\mathcal K_+(v))[N(v)]
       +D\mathcal K_+(v)[P_{<K}N(v)].
\end{align*}
In particular, no autonomous evolution is assigned to the filtered field $v$.

Taking $J=K$ leaves a quintic term whose five original velocity inputs,
nonlinear-direction output, and stretching output are all above the
same fixed cutoff. Its physical-time bound~\eqref{eq:fullyhightarget}
remains open, with the same initial-data and shorter-lifespan quantifiers
as~\eqref{eq:highrenewaltarget}. On the actual forced construction it
fails for every fixed $K,J$ and every $\theta<1$: a finite upper allowance
would bound $\mathcal G_+\geq Y/2$, contrary to the candidate's unbounded
gradient energy. This further separation pays low stretching outputs;
it does not bound regeneration among the remaining high frequencies.

\section{Verification boundary}

The Lean theorem
\texttt{openAI\_solution\_no\_uniform\_signed\_production\_budget}
applies to the actual construction at every positive viscosity and concerns
the full-field spatial cubic integral. The sharp Fourier-tail extension
and all caloric results here are analytic proofs, not Lean theorems or
axioms. The companion \texttt{labs/caloric\_force.py} checks the complete
caloric balance by exact rational Fourier arithmetic on the periodic cube,
including forces with a gradient part. It also checks the three-part
frequency split, including a nonzero low-output contribution, and the
ordered mode bound on nearly cancelling parents. Separate source-output
checks retain the original positive cost and include arbitrarily high
parent frequencies that feed low stretching outputs. These checks do not
certify the whole-space estimates or a physical-time bound.
Arbitrary-data unforced regularity remains open.

\begin{thebibliography}{9}
\bibitem{OAI}
OpenAI, \emph{Finite time blowup for Navier--Stokes},
\href{https://cdn.openai.com/pdf/32d9f210-8b73-45e0-91bc-82a30aef8a9a/navier-stokes.pdf}{manuscript},
with formalization in
\href{https://github.com/openai/NavierStokesAndEuler}{NavierStokesAndEuler}, 2026.
\end{thebibliography}
\end{document}
